\chapter{Matriks dan Graf}\label{ch:g12:matrix}

\omterm{def:g12:matrix:matrix}{Matriks} adalah larik bilangan berbentuk persegi panjang, yang dijumlahkan dan
dikalikan menurut kaidah yang dirancang agar aljabar matriksnya mewakili komposisi
transformasi linear. \omterm{def:g12:matrix:matrix}{Matriks} menyelesaikan sistem linear, menggerakkan \omterm{def:g12:seq:sequence}{barisan}
rekurensi yang berpasangan, dan mencacah jalan pada jaringan --- yaitu matematika
di balik mesin pencari dan algoritme lintasan terpendek.

\section{Aljabar matriks}

\begin{definition}[Matriks]\label{def:g12:matrix:matrix}
Suatu \emph{matriks $m \times n$}\index{matriks} adalah tabel \omterm{def:g10:numbers:sets}{bilangan real} dengan
$m$ baris dan $n$ kolom:
$A = (a_{ij})$, dengan $a_{ij}$ isian pada baris $i$, kolom $j$.
Dua matriks yang berukuran sama dijumlahkan isian demi isian, dan
$\lambda A = (\lambda a_{ij})$.
\end{definition}

\begin{definition}[Hasil kali matriks]\label{def:g12:matrix:product}
Misalkan $A$ berukuran $m \times n$ dan $B$ berukuran $n \times p$. Hasil kali $AB$
adalah \omterm{def:g12:matrix:matrix}{matriks} $m \times p$ yang isian $(i,j)$-nya
\[
(AB)_{ij} = \sum_{k=1}^{n} a_{ik} b_{kj}
\]
(yaitu kaidah ``baris $i$ dari $A$ kali kolom $j$ dari $B$'').
\end{definition}

\begin{example}
$\begin{pmatrix} 1 & 2\\ 3 & 4\end{pmatrix}
\begin{pmatrix} 0 & 1\\ 1 & 1\end{pmatrix}
= \begin{pmatrix} 2 & 3\\ 4 & 7\end{pmatrix}$,
\quad sedangkan \quad
$\begin{pmatrix} 0 & 1\\ 1 & 1\end{pmatrix}
\begin{pmatrix} 1 & 2\\ 3 & 4\end{pmatrix}
= \begin{pmatrix} 3 & 4\\ 4 & 6\end{pmatrix}$:
\emph{perkalian \omterm{def:g12:matrix:matrix}{matriks} tidak komutatif}.
\end{example}

\begin{proposition}[Kaidah aljabar matriks]\label{prop:g12:matrix:rules}
Setiap kali ukurannya membuat hasil kalinya bermakna:
\[
(AB)C = A(BC), \qquad A(B + C) = AB + AC, \qquad (A+B)C = AC + BC,
\]
dan \emph{\omterm{def:g12:matrix:matrix}{matriks} satuan} $I_n$ (berisi satu pada diagonalnya, nol di tempat lain)
memenuhi $I_m A = A I_n = A$ untuk $A$ berukuran $m \times n$.
\end{proposition}

\begin{proof}
Semuanya pemeriksaan isian demi isian dari
\cref{def:g12:matrix:product}; keasosiatifannya, satu-satunya yang tak sepele,
sama saja dengan menukar dua jumlah berhingga:
$\bigl((AB)C\bigr)_{ij} = \sum_l \left(\sum_k a_{ik}b_{kl}\right) c_{lj}
= \sum_k a_{ik} \left(\sum_l b_{kl} c_{lj}\right)
= \bigl(A(BC)\bigr)_{ij}$.
\end{proof}

\begin{definition}[Balikan]\label{def:g12:matrix:inverse}
\omterm{def:g12:matrix:matrix}{Matriks} persegi $A$ berukuran $n$ disebut \emph{dapat dibalik}\index{matriks!balikan}
jika ada \omterm{def:g12:matrix:matrix}{matriks} $B$ dengan $AB = BA = I_n$; \omterm{def:g12:matrix:matrix}{matriks} $B$ itu lalu tunggal dan ditulis
$A^{-1}$.
\end{definition}

\begin{proposition}[Balikan matriks $2\times2$]\label{prop:g12:matrix:inverse2x2}
Misalkan $A = \begin{pmatrix} a & b\\ c & d\end{pmatrix}$ dan
$\det A = ad - bc$ (yaitu \emph{determinan}\index{determinan} \omterm{def:g12:matrix:matrix}{matriks} itu). Maka $A$
\omterm{def:g12:matrix:inverse}{dapat dibalik} jika dan hanya jika $\det A \neq 0$, dan dalam hal itu
\[
A^{-1} = \frac{1}{ad - bc}\begin{pmatrix} d & -b\\ -c & a\end{pmatrix}.
\]
\end{proposition}

\begin{proof}
Perhitungan memberi
$A \begin{pmatrix} d & -b\\ -c & a\end{pmatrix}
= \begin{pmatrix} d & -b\\ -c & a\end{pmatrix} A = (ad - bc) I_2$; jika
$ad - bc \neq 0$, bagilah. Sebaliknya, jika $ad - bc = 0$, kolom $A$
sebanding, dan begitu pula kolom $AB$ untuk sebarang $B$; padahal
kolom $I_2$ tidak sebanding, jadi tak ada $B$ yang memenuhi $AB = I_2$.
\end{proof}

\begin{method}[Sistem linear]\label{met:g12:matrix:systems}
Sistem
$\begin{cases} ax + by = e\\ cx + dy = f \end{cases}$
adalah \omterm{def:g10:algebra:equation}{persamaan} \omterm{def:g12:matrix:matrix}{matriks} $AX = Y$ dengan
$X = \begin{pmatrix} x \\ y\end{pmatrix}$,
$Y = \begin{pmatrix} e \\ f\end{pmatrix}$. Jika $\det A \neq 0$, penyelesaian
tunggalnya adalah $X = A^{-1}Y$. Formalisme yang sama mengurus $n$ \omterm{def:g10:algebra:equation}{persamaan} dengan
$n$ bilangan tak diketahui.
\end{method}

\section{Pangkat matriks dan barisan rekurensi}

\begin{definition}\label{def:g12:matrix:power}
Untuk \omterm{def:g12:matrix:matrix}{matriks} persegi $A$ dan $k \in \N$, $A^k = A \times \dots \times A$
($k$ faktor), dengan $A^0 = I$.
\end{definition}

\begin{method}[Kasus diagonal-tambah-nilpoten dan kasus terdiagonalkan]\label{met:g12:matrix:powers}
Dua cara baku untuk menghitung $A^k$:
\begin{itemize}
  \item Jika $A = \lambda I + N$ dengan $N^2 = 0$, maka teorema binomial
        (yang berlaku di sini sebab $I$ dan $N$ komutatif) runtuh menjadi dua suku:
        $A^k = \lambda^k I + k \lambda^{k-1} N$.
  \item Jika orang menemukan $P$ yang \omterm{def:g12:matrix:inverse}{dapat dibalik} dengan $A = PDP^{-1}$ dan $D$
        diagonal, maka $A^k = P D^k P^{-1}$, dan $D^k$ dihitung isian demi isian.
        (Menemukan $P$ semacam itu secara sistematis adalah teori
        \emph{pendiagonalan}, yang dikembangkan di universitas; pada tingkat ini $P$
        sudah diberikan.)
\end{itemize}
\end{method}

\begin{example}[Barisan berpasangan]\label{ex:g12:matrix:coupled}
Misalkan $u_{n+1} = 3u_n + v_n$ dan $v_{n+1} = u_n + 3v_n$. Dengan mengambil
$X_n = \begin{pmatrix} u_n\\ v_n \end{pmatrix}$ dan
$A = \begin{pmatrix} 3 & 1\\ 1 & 3\end{pmatrix}$, kita peroleh
$X_{n+1} = AX_n$, sehingga $X_n = A^n X_0$. \omterm{def:g12:seq:sequence}{Barisan} bantunya
$s_n = u_n + v_n$ dan $d_n = u_n - v_n$ memenuhi $s_{n+1} = 4s_n$ dan
$d_{n+1} = 2d_n$, jadi $s_n = 4^n s_0$, $d_n = 2^n d_0$ dan
\[
u_n = \frac{4^n(u_0+v_0) + 2^n(u_0-v_0)}{2}, \qquad
v_n = \frac{4^n(u_0+v_0) - 2^n(u_0-v_0)}{2}.
\]
(Di balik layarnya: $(1,1)$ dan $(1,-1)$ adalah arah \omterm{def:g11:vect:vector}{vektor} eigen $A$.)
\end{example}

\section{Graf dan jalan}

\begin{definition}[Graf, matriks ketetanggaan]\label{def:g12:matrix:graph}
Suatu \emph{graf}\index{graf} terdiri atas simpul $1, 2, \dots, n$ dan sisi
yang menghubungkan pasangan simpul tertentu (pasangan terurut bagi graf
\emph{berarah}). \emph{Matriks ketetanggaan}\index{matriks ketetanggaan} dari graf itu
adalah \omterm{def:g12:matrix:matrix}{matriks} $n \times n$ bernama $M$ dengan $m_{ij} = 1$ jika ada sisi dari $i$ ke
$j$, dan $0$ jika tidak. Suatu \emph{jalan} sepanjang $k$ dari $i$ ke $j$ adalah
\omterm{def:g12:seq:sequence}{barisan} $k$ sisi berturut-turut yang membawa dari $i$ ke $j$.
\end{definition}

\begin{omfigure}
\begin{tikzpicture}[baseline=(current bounding box.center)]
  \node[circle, draw, inner sep=2.5pt] (v1) at (0,0) {$1$};
  \node[circle, draw, inner sep=2.5pt] (v2) at (2.6,1.2) {$2$};
  \node[circle, draw, inner sep=2.5pt] (v3) at (2.6,-1.2) {$3$};
  \draw[-Stealth, omDef, thick] (v1) -- (v2);
  \draw[-Stealth, omDef, thick] (v2) -- (v3);
  \draw[-Stealth, omDef, thick] (v1) to[bend right=18] (v3);
  \draw[-Stealth, omDef, thick] (v3) to[bend right=18] (v1);
\end{tikzpicture}
\qquad\qquad
$M = \begin{pmatrix} 0 & 1 & 1\\ 0 & 0 & 1\\ 1 & 0 & 0 \end{pmatrix}$

{\small Sebuah \omterm{def:g12:matrix:graph}{graf} berarah dan \omterm{def:g12:matrix:graph}{matriks ketetanggaannya}
(\cref{exo:g12:matrix:6}): $m_{ij} = 1$ tepat ketika ada sisi
dari $i$ ke $j$.}
\end{omfigure}

\begin{theorem}[Mencacah jalan]\label{thm:g12:matrix:walks}
Banyaknya jalan sepanjang $k$ dari simpul $i$ ke simpul $j$ adalah
isian $(i,j)$ pada $M^k$.
\end{theorem}

\begin{proof}
Induksi pada $k$. Untuk $k = 1$ inilah definisi $M$. Andaikan
pernyataannya benar untuk $k$. Jalan sepanjang $k+1$ dari $i$ ke $j$ adalah jalan
sepanjang $k$ dari $i$ ke suatu simpul $l$, lalu diikuti sisi dari $l$ ke $j$; menurut
asas penjumlahan dan asas perkalian, banyaknya adalah
\[
\sum_{l=1}^{n} \bigl(M^k\bigr)_{il}\, m_{lj} = \bigl(M^{k+1}\bigr)_{ij}.
\qedhere
\]
\end{proof}

\begin{example}
Untuk \omterm{def:g12:matrix:graph}{graf} segitiga ($3$ simpul, semua pasangannya terhubung),
$M = \begin{pmatrix} 0&1&1\\ 1&0&1\\ 1&1&0\end{pmatrix}$ dan
$M^2 = \begin{pmatrix} 2&1&1\\ 1&2&1\\ 1&1&2\end{pmatrix}$: jadi dari tiap simpulnya
ada $2$ jalan sepanjang $2$ yang kembali ke dirinya (lewat salah satu tetangganya) dan
$1$ jalan ke tiap simpul lainnya.
\end{example}

\section{Latihan}

\begin{exercise}[$\star$]\label{exo:g12:matrix:1}
Misalkan $A = \begin{pmatrix} 1 & 2\\ 0 & 1 \end{pmatrix}$ dan
$B = \begin{pmatrix} 2 & 0\\ 1 & 1 \end{pmatrix}$. Hitunglah $A + B$, $AB$,
$BA$ dan $A^2$.
\end{exercise}

\begin{exercise}[$\star$]\label{exo:g12:matrix:2}
Tentukan apakah \omterm{def:g12:matrix:matrix}{matriks} berikut \omterm{def:g12:matrix:inverse}{dapat dibalik}, lalu hitunglah
balikannya jika ada:
\[
A = \begin{pmatrix} 2 & 5\\ 1 & 3\end{pmatrix}, \qquad
B = \begin{pmatrix} 3 & 6\\ 2 & 4\end{pmatrix}.
\]
\end{exercise}

\begin{exercise}[$\star$]\label{exo:g12:matrix:3}
Selesaikan lewat pembalikan \omterm{def:g12:matrix:matrix}{matriks} sistem
$\begin{cases} 2x + 5y = 1\\ x + 3y = 2 . \end{cases}$
\end{exercise}

\begin{exercise}[$\star\star$]\label{exo:g12:matrix:4}
Misalkan $A = \begin{pmatrix} 2 & 1\\ 0 & 2\end{pmatrix} = 2I + N$ dengan
$N = \begin{pmatrix} 0 & 1\\ 0 & 0 \end{pmatrix}$.
\begin{enumerate}
  \item Periksalah bahwa $N^2 = 0$ dan bahwa $I$ dan $N$ komutatif.
  \item Simpulkan $A^k$ untuk semua $k \in \N$ lalu periksalah rumusnya untuk $k=2$
        lewat perhitungan langsung.
\end{enumerate}
\end{exercise}

\begin{exercise}[$\star\star$]\label{exo:g12:matrix:5}
Misalkan $A = \begin{pmatrix} 0 & 1\\ 1 & 1\end{pmatrix}$ dan $F_n$ \omterm{def:g12:seq:sequence}{barisan}
Fibonacci ($F_0 = 0$, $F_1 = 1$, $F_{n+2} = F_{n+1} + F_n$). Tunjukkan lewat
induksi bahwa untuk $n \geq 1$,
\[
A^n = \begin{pmatrix} F_{n-1} & F_n\\ F_n & F_{n+1}\end{pmatrix},
\]
lalu simpulkan identitas $F_{n+1}F_{n-1} - F_n^2 = (-1)^n$.
(Petunjuk: \omterm{def:g11:vect:det}{determinan} saling mengalikan: $\det(MN) = \det M \det N$, dan itu boleh
kamu periksa untuk \omterm{def:g12:matrix:matrix}{matriks} $2\times2$.)
\end{exercise}

\begin{exercise}[$\star\star$]\label{exo:g12:matrix:6}
Sebuah \omterm{def:g12:matrix:graph}{graf} berarah pada simpul $\{1, 2, 3\}$ mempunyai sisi
$1\to2$, $2\to3$, $3\to1$ dan $1\to3$.
\begin{enumerate}
  \item Tulislah \omterm{def:g12:matrix:graph}{matriks ketetanggaannya} $M$ lalu hitung $M^2$ dan $M^3$.
  \item Ada berapa jalan sepanjang $3$ dari $1$ ke $1$? Daftarkanlah.
\end{enumerate}
\end{exercise}

\begin{exercise}[$\star\star$]\label{exo:g12:matrix:7}
Sebuah perusahaan berbagi mobil memindahkan kendaraan antara dua kota $A$ dan $B$.
Tiap pekan, $80\%$ mobil di $A$ tinggal di $A$ dan $20\%$ pindah ke $B$; lalu $30\%$
mobil di $B$ pindah ke $A$ dan $70\%$ tinggal. Misalkan $a_n, b_n$ adalah
bagian armadanya di tiap kota.
\begin{omfigure}
\begin{tikzpicture}
  \node[circle, draw, inner sep=3pt] (A) at (0,0) {$A$};
  \node[circle, draw, inner sep=3pt] (B) at (3.2,0) {$B$};
  \draw[-Stealth, omDef, thick] (A) to[bend left=25]
    node[above, font=\small] {$0.2$} (B);
  \draw[-Stealth, omDef, thick] (B) to[bend left=25]
    node[below, font=\small] {$0.3$} (A);
  \draw[-Stealth, omDef, thick] (A) to[loop left]
    node[left, font=\small] {$0.8$} (A);
  \draw[-Stealth, omDef, thick] (B) to[loop right]
    node[right, font=\small] {$0.7$} (B);
\end{tikzpicture}
\end{omfigure}
\begin{enumerate}
  \item Tulislah $X_{n+1} = MX_n$ dengan
        $X_n = \begin{pmatrix} a_n\\ b_n\end{pmatrix}$ lalu kenalilah $M$.
  \item Carilah bagian yang setimbang (selesaikan $MX = X$ dengan
        $a + b = 1$).
  \item Tunjukkan bahwa $c_n = a_n - 0.6$ memenuhi $c_{n+1} = 0.5\,c_n$, lalu
        simpulkan bahwa \omterm{def:g11:prob:rv}{distribusi} armadanya menuju kesetimbangan itu.
\end{enumerate}
\end{exercise}

\begin{exercise}[$\star\star\star$]\label{exo:g12:matrix:8}
Misalkan $A = \begin{pmatrix} 3 & 1\\ 1 & 3 \end{pmatrix}$,
$P = \begin{pmatrix} 1 & 1\\ 1 & -1 \end{pmatrix}$.
\begin{enumerate}
  \item Hitunglah $P^{-1}$, lalu $D = P^{-1}AP$, dan periksalah bahwa $D$
        diagonal.
  \item Simpulkan rumus tertutup bagi $A^n$ lalu bandingkan dengan
        \cref{ex:g12:matrix:coupled}.
\end{enumerate}
\end{exercise}

\section{Soal: Matriks yang hafal Fibonacci (dan
cuaca)}

\begin{problem}[{Soal akhir pekan --- satu matriks $2 \times 2$
memikul seluruh Fibonacci, sebuah matriks Markov meramal cuaca jangka
panjang, dan sebuah vektor eigen bernilai semiliar dolar}]
\label{pb:g12:matrix:1}
\omterm{def:g12:matrix:matrix}{Matriks} adalah mesin yang memakan sebuah keadaan lalu mengembalikan keadaan
berikutnya --- dan \emph{pangkatnya} karena itu menyimpan seluruh masa depan. Soal
ini dibuka dengan \omterm{def:g12:matrix:matrix}{matriks} mencengangkan yang pangkatnya mendaftar
bilangan Fibonacci (dan membuktikan identitasnya masing-masing satu baris),
lalu menjalankan cuaca sebagai rantai Markov sampai keadaan mapannya, dan
ditutup dengan \omterm{def:g11:vect:vector}{vektor} eigen yang di atasnya sebuah mesin pencari dibangun
(\cref{thm:g12:matrix:walks},
\cref{met:g12:matrix:powers}).

\medskip
\noindent\textbf{Bagian I --- Kelancaran.}
\begin{enumerate}
  \item Dengan $A = \begin{pmatrix} 1 & 2\\ 3 & 4\end{pmatrix}$
        dan $B = \begin{pmatrix} 0 & 1\\ 1 & 0\end{pmatrix}$:
        hitunglah $AB$ dan $BA$. Vonis atas kekomutatifannya?
  \item Baliklah $\begin{pmatrix} 2 & 1\\ 5 & 3\end{pmatrix}$
        (\cref{prop:g12:matrix:inverse2x2}) lalu pakailah balikannya
        untuk menyelesaikan $2x + y = 4$, $5x + 3y = 7$.
  \item Misalkan $N = \begin{pmatrix} 0 & 1\\ 0 & 0\end{pmatrix}$:
        hitunglah $N^2$, lalu simpulkan
        $(I + N)^n = I + nN$ untuk setiap $n$.
  \item \omterm{def:g12:matrix:graph}{Graf} segitiga (tiga simpul, semua pasangannya terhubung):
        tulislah \omterm{def:g12:matrix:graph}{matriks ketetanggaannya} $A$, hitunglah $A^3$, lalu
        tafsirkan isian diagonalnya
        (\cref{thm:g12:matrix:walks}).
  \item Untuk $D = \begin{pmatrix} 2 & 0\\ 0 & \frac12
        \end{pmatrix}$: berikan $D^n$ dan perilakunya ketika
        $n \to \infty$.
\end{enumerate}

\medskip
\noindent\textbf{Bagian II --- \omterm{def:g12:matrix:matrix}{Matriks} Fibonacci.}
Misalkan $F = \begin{pmatrix} 1 & 1\\ 1 & 0\end{pmatrix}$ dan misalkan
$F_1 = F_2 = 1, F_3 = 2, \dots$ adalah bilangan Fibonacci pada
\cref{pb:g11:seq:1}.
\begin{enumerate}[resume]
  \item Hitunglah $F^2$, $F^3$, $F^4$ lalu terkalah bentuk umum
        $F^n$ dengan bilangan Fibonacci.
  \item Buktikan terkaan
        $F^n = \begin{pmatrix} F_{n+1} & F_n\\ F_n & F_{n-1}
        \end{pmatrix}$ itu lewat induksi.
  \item Ambillah \omterm{def:g11:vect:det}{determinan} kedua ruasnya (\omterm{def:g11:vect:det}{determinan} hasil
        kali adalah hasil kali \omterm{def:g11:vect:det}{determinannya} --- periksalah
        pada \omterm{def:g12:matrix:matrix}{matriks} $2 \times 2$ jika kamu belum pernah melihatnya):
        lalu simpulkan \emph{identitas Cassini}
        $F_{n+1}F_{n-1} - F_n^2 = (-1)^n$ --- yaitu mesin
        kuadrat yang lenyap itu, yang terbukti dalam satu baris.
  \item Dari $F^{m+n} = F^m F^n$, bacalah isian
        kanan atasnya lalu turunkan \emph{rumus penjumlahannya}
        \[
        F_{m+n} = F_{m+1} F_n + F_m F_{n-1} .
        \]
        Periksalah untuk $m = n = 3$.
  \item Simpulkan dari rumus penjumlahan itu (lewat induksi pada $k$)
        bahwa $F_n$ \omterm{def:g12:arith:divides}{membagi} $F_{kn}$, lalu periksalah pada
        $F_3 \mid F_6$ dan $F_3 \mid F_9$.
  \item Untuk menghitung $F_{100}$, orang tak perlu mengalikan $100$
        \omterm{def:g12:matrix:matrix}{matriks}: kuadratkanlah berulang kali ($F^2, F^4, F^8, \dots$)
        lalu gabungkan. Berapa kali perkalian \omterm{def:g12:matrix:matrix}{matriks} yang cukup,
        dan muslihat perkalian kuno mana dari jilid sekolah
        menengah pertama yang ini, yang kini naik pangkat ke \omterm{def:g12:matrix:matrix}{matriks}?
\end{enumerate}

\medskip
\noindent\textbf{Bagian III --- Mesin cuaca.}
Di sebuah kota: sesudah hari yang cerah, hari berikutnya cerah dengan
peluang $0.8$; sesudah hari yang hujan, cerah dengan peluang
$0.4$. Sandikan \omterm{def:g11:prob:rv}{distribusi} harinya sebagai satu kolom
$\binom{p_{\text{cerah}}}{p_{\text{hujan}}}$ dan perubahannya oleh
\[
M = \begin{pmatrix} 0.8 & 0.4\\ 0.2 & 0.6 \end{pmatrix}.
\]
\begin{enumerate}[resume]
  \item Periksalah bahwa tiap kolom $M$ berjumlah $1$, lalu katakan mengapa
        setiap mesin cuaca harus punya sifat itu.
  \item Hari ini cerah. Hitunglah ramalan untuk besok dan
        untuk lusa.
  \item Carilah \emph{keadaan mapannya}: yaitu \omterm{def:g11:prob:rv}{distribusi} $v$ dengan
        $Mv = v$ (dan isiannya berjumlah $1$). Berapa bagian
        harinya yang cerah dalam jangka panjang?
  \item Mulailah dari hari hujan, $\binom01$, lalu terapkan $M$ empat
        kali, sambil melacak jaraknya ke keadaan mapan pada
        tiap langkahnya. Dengan faktor berapa jurangnya menyusut tiap langkah
        --- dan kekonvergenan macam apa ini?
  \item PageRank dalam ukuran mini: tiga halaman, dengan tautan
        $A \to B$, $A \to C$, $B \to C$, $C \to A$. Seorang peselancar
        acak mengikuti tautan keluar secara acak seragam.
        Tulislah \omterm{def:g12:matrix:matrix}{matriks} peralihannya, carilah keadaan mapannya, lalu
        peringkatkan halamannya.
  \item Tafsirkan peringkatnya: mengapa $C$ bernilai setinggi
        $A$ padahal menerima tautan dari lebih sedikit halaman --- apa
        yang sebenarnya diukur keadaan mapan itu? (PageRank yang sungguhan
        menambahkan faktor peredam untuk jalan buntu dan lompatan; gagasan
        \omterm{def:g11:vect:vector}{vektor} eigennya tepat yang ini.)
\end{enumerate}

\medskip
\noindent\textbf{Bagian IV --- Untung dari diagonalnya.}
\begin{enumerate}[resume]
  \item Dua besaran yang berpasangan menuruti
        $u_{n+1} = 3u_n + v_n$, $v_{n+1} = u_n + 3v_n$, yaitu
        \omterm{def:g12:matrix:matrix}{matriks} $A$ pada \cref{exo:g12:matrix:8}. Dengan memakai
        pendiagonalan latihan itu
        ($D = \operatorname{diag}(4, 2)$), berikan rumus
        tertutup bagi $u_n$ ketika $u_0 = 1$, $v_0 = 0$, lalu periksalah
        terhadap perhitungan langsung untuk $n = 1, 2, 3$.
  \item Dalam satu atau dua kalimat: apa yang \emph{dilakukan} pendiagonalan
        terhadap sistem yang berpasangan --- dan dalam arti apa
        keadaan mapan Markov pada pertanyaan 14 juga sebuah kisah \omterm{def:g11:vect:vector}{vektor}
        eigen?
  \item Penutup --- tiga wajah \omterm{def:g12:matrix:matrix}{matriks} pada akhir pekan ini:
        pembukuan (sistem dan balikan), kombinatorika
        (jalan dan tautan yang tercacah oleh pangkatnya), dan perubahan
        (Fibonacci, cuaca, jejaring --- masa depan yang terbaca dari
        arah eigennya). Masing-masing satu kalimat, ditambah
        penunjuk ke depan: aljabar linear di jilid universitas
        menjadikan tiap wajah itu sebuah teori.

\end{enumerate}
\end{problem}
